\item \textbf{Visualización de la distribución estelar.}

\begin{tcolorbox}[colback=blue!2!white,colframe=blue!55!black,title={Enunciado}]

A partir de los datos generados en el problema anterior, se desea visualizar la distribución de estrellas en la galaxia.

\begin{itemize}
\item[(a)] Use \texttt{matplotlib.pyplot} para construir un gráfico de dispersión (\texttt{scatter}) de las posiciones \((x,y)\).

\item[(b)] Configure el gráfico de manera que:
\begin{itemize}
\item tenga título,
\item tenga etiquetas en los ejes,
\item incluya una grilla.
\end{itemize}

\item[(c)] Utilice la distancia al origen \(r_i\) para asignar un color a cada punto en el gráfico (usando el argumento \texttt{c=} en \texttt{scatter}).

\item[(d)] Añada una barra de color (\texttt{colorbar}) para visualizar cómo varía la distancia al centro.

\end{itemize}

\end{tcolorbox}

\begin{solution}

\subsubsection*{(a) Gráfico de dispersión}

Se utiliza \texttt{matplotlib.pyplot} para visualizar las posiciones:

\begin{pythonjupyter}
import matplotlib.pyplot as plt

plt.scatter(x, y)
plt.show()
\end{pythonjupyter}

\subsubsection*{(b) Configuración del gráfico}

Se añaden elementos básicos para mejorar la visualización:

\begin{pythonjupyter}
plt.figure(figsize=(6,6))
plt.scatter(x, y)

plt.title("Distribución de estrellas en la galaxia")
plt.xlabel("Posición en x")
plt.ylabel("Posición en y")
plt.grid()

plt.show()
\end{pythonjupyter}

\subsubsection*{(c) Uso del color según la distancia}

Se utiliza el arreglo \texttt{r} para asignar colores a cada punto:

\begin{pythonjupyter}
plt.figure(figsize=(6,6))
plt.scatter(x, y, c=r)

plt.title("Distribución de estrellas (color según distancia)")
plt.xlabel("Posición en x")
plt.ylabel("Posición en y")
plt.grid()

plt.show()
\end{pythonjupyter}

\subsubsection*{(d) Barra de color}

Se agrega una barra de color para interpretar los valores:

\begin{pythonjupyter}
plt.figure(figsize=(6,6))
scatter = plt.scatter(x, y, c=r)

plt.title("Distribución de estrellas (color según distancia)")
plt.xlabel("Posición en x")
plt.ylabel("Posición en y")
plt.grid()

plt.colorbar(scatter, label="Distancia al origen")

plt.show()
\end{pythonjupyter}

El resultado permite visualizar cómo las estrellas se distribuyen alrededor del centro y cómo varía su distancia mediante el color.

\end{solution}
